\section{Limitations and Trust Assumptions}
\label{sec:limitations}

This section lists what the design does not protect against and the trust assumptions it makes.

\begin{caveatbox}
The dead-man-switch protects against the owner going \emph{silent}, not against a
compromised key while the owner is active: anyone holding sufficient owner
authority can spend and reset the timer. Liveness also depends on the off-chain
wallet renewing proof of life on the owner's behalf. The on-chain timer
enforces the deadline, but choosing a sensible timeframe and trustworthy
beneficiaries remains the owner's responsibility.
\end{caveatbox}

\begin{description}[leftmargin=1.4em,style=nextline]
\item[Broad operator spend (by design, G3).] A quorum of owner keys can drain the wallet down to the streaming reserve. The mitigations are the dead-man-switch, the recovery-reachability invariant, multi-signature thresholds, and trusted transaction construction, not on-chain destination or amount limits. Narrowing the operator envelope would be a product change, not a security fix.
% VERIFIED: lib/state/configuration.ak:140-143 and lib/stt/operator_handlers.ak:103-128 permit independent grants through authorized configuration.
\item[Allowance grants and operator changes.] Configuration validates the remaining and per-day allowances independently.
The initial remaining grant may exceed the per-day grant for an asset.
An authorized operator may also refill or reconfigure a user's allowance through $\mathrm{UpdateState}$.
These are accepted grants, not unauthorized withdrawals.
The earlier unconditional daily-bound claim omitted these cases.
Theorem~\ref{thm:allowance-velocity} now states the required starting bound and excludes operator changes during the considered interval.
\item[Wallet-less allowance succession.] A signed allowance user may submit a non-empty \texttt{UseAllowance} transition without consuming a wallet input. The STT-only form reduces the user's remaining allowance but moves no wallet value. It may repeat against each recreated STT until no positive effective allowance remains. A later reset supplies the next configured grant. Each call needs external fee funding and has no cadence rule.
% VERIFIED: lib/stt/user_handlers.ak::eval_use_allowance has no wallet-input requirement. validators/stt_allowance_tests.ak::allowance_use_accepts_wallet_less_draw_from_valid_configuration executes the STT-only form.
% INFERRED: each confirmed successor invalidates a pending transaction that references the prior STT output.
\item[Positive-power records require distinct credentials.] At mint and full state update, each payment key hash may appear in at most one user record with positive multi-signature power. Each signed record contributes its configured weight once. Records with no power or zero power may share credentials, and a user credential may also identify a beneficiary. Distinct credentials do not prove distinct people. One person who controls several keys can still activate several weighted records. The frontend applies the same rule before transaction construction.
% VERIFIED: code/smart-contract/lib/state/configuration.ak::expect_users_are_valid; validators/stt_mint_tests.ak::stt_mint_rejects_shared_multisig_credential; validators/stt_operator_tests.ak::operator_state_update_rejects_shared_multisig_credential; code/dApp/src/lib/contracts/state-validation.test.ts::rejects one credential across multiple positive-power owner records.
\item[Admin-only wallets are permitted.] A wallet may be configured with a sole admin and no other recovery path; losing that key is then an accepted user risk, exactly as for any single-signer wallet. The reachability invariant only guarantees that an \emph{admin-less} wallet always retains a non-admin way back in.
\item[Shape, not timing, of recovery.] The contract rejects only genuinely bricked configurations; a beneficiary whose unlock lies far in the future is accepted on-chain and flagged off-chain. A trusted liveness keeper can defer beneficiary unlock indefinitely while it keeps renewing, and an operator (admin or quorum) can push the deadline far forward in a single reconfiguration: the one-increment cap binds only the heartbeat renewal paths, not the full state-update path, which must be able to (re)configure the timer from scratch.
\item[Authorized proof-of-life renewal contention.] The exact no-op guard does not set a minimum deadline increase or renewal cadence. A keeper can advance the deadline by one unit per accepted successor while each value remains inside the renewal window. Each successor can invalidate pending operator or beneficiary transactions that use the prior state token. The keeper pays every fee and cannot move wallet value through this action. The 2026-09 security review accepts this availability risk because the keeper is a trusted liveness role.
% VERIFIED: lib/stt/user_handlers.ak::eval_renew_proof_of_life requires a changed deadline. lib/state/proof_of_life.ak::expect_valid_renewal_window has no minimum increase or cadence. The user accepted this residual during the 2026-09 security review.
% INFERRED: a confirmed successor invalidates a competing transaction that uses the prior STT.
\item[Streaming payments are additive and end-date governed.] An operator (admin or quorum) may add streaming payments to a live wallet or configure them at mint. Each addition is unsettled, and the count cap bounds the set. Management cannot delete an existing payment or change its immutable fields. It can move the end later, or earlier down to the transaction's upper validity bound. This rule protects already-accrued value. A payment's own payee may shorten its stream through the same no-clawback rule. Before the start, it may cap the end at the start and create a zero-lifetime obligation. There is no cancellation marker. Before terminal recovery, the payee may shorten it again after the shared cooldown. During terminal recovery, the payee must use the exact safe cutoff and cannot repeat for that payment. An operator may later reschedule it. A payment leaves the set only through settlement. The contract recomputes the per-asset reserve from the live payment set on every spend.
\item[STT-policy streaming assets.] On-chain stream ingress rejects the STT script as a payout payment credential, but it does not reject the STT policy as the payment asset policy. If such a stream ever owes a positive integer quantity, it cannot make positive payment progress or leave State while that unpaid amount remains. Settlement cannot put that policy's asset at its tagged non-STT payee. Each STT token must remain at its own continuing STT output, and the spend path accepts only one input at the shared STT address. Final permanent exit and exact distribution require an empty stream list. Reserve-aware beneficiary recovery can still recover other wallet assets. The maintained dApp rejects the STT policy for fresh mint and management entries. It permits existing entries to remain forwardable and reschedulable. Custom builders must apply the same exclusion.
% VERIFIED: lib/streaming_payments/shape.ak checks asset identity but has no STT-policy exclusion. validators/stt_operator_tests.ak::manage_streaming_payments_accepts_new_stream_asset_under_stt_policy pins the on-chain boundary. The consumed STT supplies asset evidence. Correction: the former mint acceptance test relied only on the newly minted STT. validators/stt_mint_tests.ak::stt_mint_rejects_stream_asset_only_in_current_stt_mint now rejects that transaction because the asset is absent from both input lists. code/dApp/src/lib/contracts/state-validation.test.ts, streaming-manage.test.ts, and lib/mesh/transactions/mint-state-token.test.tsx cover the maintained dApp guard.
% INFERRED: the positive-payout deadlock follows from lib/streaming_payments/payout.ak, lib/stt/io.ak, and the empty-stream gates in lib/stt/user_handlers.ak.
% VERIFIED: code/smart-contract/lib/streaming_payments/asset_presence.ak::new_assets_are_present and validators/stt_streaming_asset_presence_tests.ak.
\item[Streaming asset existence.] Input evidence proves that an asset exists when the payment is added. The input can belong to any address, including another holder's reference input. This check does not prove wallet ownership, sufficient payment funding, or future asset availability. One base unit can supply evidence for several new payments of the same asset. Newly minted tokens and transaction outputs alone cannot supply evidence. Existing forwarded or rescheduled payments need no new evidence. A payment asset must therefore exist before the transaction that introduces it as a stream.
\item[Streaming reserve is point-in-time.] The reserve protects only streaming-payment value \emph{accrued} as of a spend's upper validity bound, not the full remaining schedule. A beneficiary recovering after the owner goes silent can therefore draw down principal that a payment's \emph{future} accrual will need, leaving later settlements underfunded. Reserving the whole forward schedule would let long-dated payments starve heirs, so the design protects only accrued value; keeping enough on hand for future accrual remains the owner's planning responsibility.
\item[Permanent state thread.] There is no burn/close path; the small state-token deposit persists even after funds are fully recovered. This keeps the ``exactly one STT, always'' invariant unconditional.
\item[Open settlement timing.] While a transaction's lower bound is before terminal recovery, any listed user, stream payee, or unlocked beneficiary may trigger a streaming payout and pay its fee. An admin may also settle. Once that lower bound reaches the final beneficiary's unlock boundary, only an admin or that beneficiary may trigger a payout. The destination and payout amount remain fixed by State. Before terminal recovery, cadence limits how often other stakeholders may settle, but the owner does not control when they choose to act. Parties with no relationship to the wallet cannot settle.
\item[Under-funded settlement is first-come, not pro-rata.] The reserve gate is a floor on discretionary outflow, and settlement is exempt from it: a settlement's outflow is independently pinned, asset by asset, to the tagged payee outputs, so it cannot underfund anyone and needs no floor. Applying the floor to settlement was in fact a deadlock --- when a wallet holds less of an asset than it owes, the floor forbids \emph{every} outflow of that asset, including the only action that can reduce the debt, so the asset would freeze permanently for payee, operator and beneficiary alike while accrual kept growing. The exemption removes the freeze at the cost of ordering: when a wallet cannot cover everything it owes, payees are settled in the order they crank rather than proportionally. Each payment is still individually capped at its own accrual, so no payee is ever paid more than it earned; keeping the wallet funded above its accrued obligations remains the owner's planning responsibility.
\item[Repeatable final-beneficiary recovery.] In the existing weighted withdrawal path, every beneficiary before the final one is removed after one recovery action. Removing the native-asset count cap does not let an earlier beneficiary retry an incomplete withdrawal. Ledger limits may still prevent the selected wallet inputs from fitting one transaction. The final beneficiary stays in State. It can recover separate current or future wallet UTxOs through repeated transactions. Each value-moving action may select any wallet-input set that fits ledger limits. It may leave continuing outputs for the streaming reserve or a smaller chosen withdrawal. No transaction can prove that another wallet UTxO does not exist or that no future deposit will arrive. The contract therefore has no final recovery marker or atomic all-pool guarantee. Each attempt obeys the shared thirty-minute cadence, the streaming reserve, and ledger resource limits. This path recovers wallet UTxOs only. Staking-reward withdrawal remains operator-only.
\item[Uncadenced exact-distribution succession.] With two or more beneficiaries, exact distribution requires an empty stream list and every beneficiary to be unlocked. The named initiating beneficiary must sign. A wallet-less form preserves State, including the cadence stamp. After acceptance, the signer can immediately build a successor against the recreated STT. Each successor spends and recreates the latest STT, so this is not byte replay. The caller supplies the fee, and no wallet value moves. A successor can invalidate another pending transaction that references the prior STT. This uncadenced succession and its STT contention are accepted design trade-offs. The sole-beneficiary form remains subject to the shared cooldown and window cap.
% VERIFIED: code/smart-contract/lib/stt/user_handlers.ak::eval_distribute_beneficiaries preserves State and multi-beneficiary cadence; validators/beneficiary_distribution_tests.ak::authorized_stt_only_transaction_preserves_multiple_beneficiary_state executes the wallet-less form.
% INFERRED: Singleton STT succession makes the accepted call conflict with a pending transaction that references the prior output.
% VERIFIED: exact distribution rejects nonempty streams and fractional shares.
% INFERRED: funding, lost-key, discovery, and transaction-size limitations follow from those gates and the preserved operator paths.
\item[Exact distribution requires settled streams and divisible assets.] One underfunded stream blocks exact distribution even when it owes another asset. Additional funding or the existing reserve-aware weighted withdrawal path remains necessary. An indivisible asset cannot be split among several weighted beneficiaries. Permanent exits can leave one beneficiary, but lost keys can prevent that outcome. Exact distribution preserves rights to unselected and future deposits. It does not freeze operators, guarantee discovery of all funds, or guarantee that every one-input payout fits the ledger limits.
\item[Authorized Consolidation controls wallet layout.] An admin, quorum, or unlocked beneficiary may preserve aggregate wallet Value while increasing the number of continuing wallet UTxOs. The ledger requires minimum ada in each output and enforces transaction byte-size and ExUnit limits. A later authorized Consolidation can reverse the split, but discovery work and cleanup fees can increase.
\item[Value held at the STT output.] Every otherwise-valid transition may add externally funded ada to the continuing STT output. Only an admin or a multi-signature quorum using \texttt{RunOperator(Use)} may remove ada or add or remove assets under other policies. The current wallet STT itself remains fixed. A beneficiary-only State cannot recover a top-up from this output, so the top-up remains locked unless usable operator authority exists.
% VERIFIED: lib/stt/io.ak::stt_value_preserved_with_lovelace_top_up, lib/stt/operator_handlers.ak::eval_operator_use, and validators/stt_spend_value_tests.ak.
% CORRECTION: earlier wording incorrectly excluded third-party discovery griefing because the sender funds the dust. The loss and the discovery burden affect different parties.
% VERIFIED: stt/io.ak::expect_single_stt_io and dApp transactions/internals/state-forwarding.ts::resolveStateForwardingInput.
\item[Only token-bearing outputs at the state-token address are spendable.] Every state token under the policy shares one address. A UTxO sent there without a state token fails the token check when spent alone. Spending it alongside the real state token fails the single-input-at-this-address check. Anyone can send such outputs to the public address. Their value is inaccessible under this validator. Accumulated outputs can burden a client that enumerates the address. The reference builder fetches the chosen transaction's outputs and checks the exact State token. This keeps unrelated address outputs out of that lookup.
\item[Self-addressed streams.] A stream may name the wallet's own full address on-chain. Mint and stream addition exclude the state-token payment credential but do not exclude the wallet payment credential. The maintained dApp rejects this credential at mint and for new stream additions. Existing entries remain payable and removable. A custom builder can still submit the accepted on-chain configuration. When settlement consumes wallet UTxOs, a tagged output at the wallet address counts as both the configured payout and a continuing wallet output. For ADA, the continuing wallet aggregate must equal the wallet input aggregate minus the validated payout delta. ADA routing permits ADA at all correctly tagged configured stream outputs and checks only the aggregate tagged ADA amount. The exact wallet outflow can therefore reach another configured stream payee. The validators do not bind individual lovelace from one stream's delta to that stream's address. A wallet-less settlement can use external value to create a tagged wallet output because the wallet validator does not run.

For a native asset, exact tagged-quantity matching still applies. It does not prove external delivery when the wallet is the configured payee. The acceptance test starts with two wallet inputs holding 100 tokens. It creates a tagged 10-token wallet output, an untagged 80-token wallet output, and a 10-token burn. The wallet aggregate falls from 100 to 90 while the tagged quantity remains equal to the payout delta. The Epora validators accept this transaction context. The test uses a zero fee and does not execute a minting policy or Cardano phase-one validation. A ledger-balanced form of this native-only fixture needs external ADA for a positive fee and a minting policy that accepts the burn. No external payee receives the burned quantity. Total wallet loss remains capped by the validated payout delta. These self-addressed outcomes are accepted configuration risks.
% VERIFIED: validators/stt_mint_tests.ak::stt_mint_accepts_stream_to_wallet_address; validators/stt_operator_tests.ak::manage_streaming_payments_accepts_adding_wallet_address_stream; validators/wallet_rule_tests.ak::streaming_payment_payout_rule_accepts_configured_wallet_self_address; validators/wallet_spend_tests.ak::streaming_payment_payout_accepts_self_address_delta_at_other_payee; validators/wallet_spend_tests.ak::streaming_payment_payout_accepts_self_address_native_burn. The ADA test runs both validators with no external input and sends a 1 ADA wallet delta to another configured stream payee. The native test runs the STT validator and both wallet input legs, leaves 90 of 100 tokens in the wallet, and burns 10.
% VERIFIED: code/dApp/src/lib/mesh/transactions/mint-state-token.test.tsx rejects the derived wallet destination at mint. code/dApp/src/lib/contracts/streaming-manage.test.ts rejects fresh wallet-credential destinations across stake variants and accepts an existing entry.
% VERIFIED: The native acceptance test uses the fixture's zero placeholder fee and does not execute a minting policy or Cardano phase-one validation.
% INFERRED: A ledger-balanced form of this native-only fixture needs external ADA for a positive fee and a minting policy that accepts the negative mint.
% INFERRED: Ledger value conservation can assign part or all of the exact ADA wallet delta to transaction fees when outputs do not consume it.
\item[External script payout compatibility.] A streaming or exact beneficiary payout may target an unrelated script address. The payout carries the protocol's inline $\mathrm{OutputId}$ datum. Epora validates the complete address, tag, and amount. It cannot validate whether the receiving script later accepts that datum. The person who configures the payout must verify recipient compatibility. Another Epora wallet is compatible because its spending validator ignores input datums.
% VERIFIED: lib/streaming_payments/payout.ak::output_with_streaming_payment_id_matches and lib/wallet/beneficiary_distribution.ak::all_shares_are_paid require the configured address and OutputId. lib/wallet/spend.ak::eval_spend ignores its input datum.
\item[ADA settlement granularity and fee funding.] The payout-integrity routing (Section~\ref{sec:security}) forbids \emph{any} non-protocol output from carrying the payout asset, and every Cardano output must carry a minimum amount of ADA. For an ADA-denominated stream the payout asset \emph{is} ADA, with two consequences. First, a settlement funded from wallet value cannot emit the cranker's own fee-change output --- that untagged change output would carry the payout asset and be rejected --- so the cranker's funding input has only two legal ADA sinks, the tagged payee output and the transaction fee, and is allocated across the two with no residual returned (the fee is \emph{not} equal to the input once a top-up is present: part of the input tops the payee output up to the min-UTxO amount described next, and the remainder pays the fee). Second, the settled amount can be smaller than the per-output ADA minimum, and an exact-delta payee output would then be unconstructible. To keep small ADA settlements possible, the tagged-output check for ADA requires the payee outputs to sum to \emph{at least} the delta rather than exactly it (Section~\ref{sec:security}, ``Payout integrity''). For a payout address that does not share the wallet payment credential, the cranker tops the payee's tagged output up to a valid amount with its own ADA. The wallet's net outflow remains pinned to exactly the delta and can reach only configured payees. In that non-self-addressed case, any surplus above the delta comes from the cranker. The self-addressed case above is the source-attribution exception. Its tagged wallet continuation can satisfy the stream floor while the exact wallet delta reaches another configured stream output. The residual cost is per-payee exactness across simultaneous ADA streams (see ``Payout integrity''). Constructing the crank (no change output; sufficient top-up on small ADA settlements) is the off-chain builder's responsibility. A non-ADA stream normally avoids this minimum-output issue because its padding uses ADA. The self-address exception also applies under the inferred ledger conditions above. A tagged wallet continuation can match the exact delta while the burn supplies the net wallet outflow. The self-address limitation documents this accepted risk.
\end{description}
